% !TEX root = ../nc-quadric-classification.tex
% 2026-11-22 / v1.9（全証明はBへ移管、主張のみ保持）
\section{固定した相対切断環の二次 AS 被覆}
\label{sec:relative-as-cover}

平方根を固定した相対軌道環が既に得られている場合には、
その AS 被覆の存在を有限表示の問題へ帰着できる。
以下では被覆を選ぶ空間と、被覆の存在だけを要求する部分関手を区別する。
特に前者の代数性だけから後者の代数性を結論しない。

\begin{theorem}[二次被覆の相対有限判定]
\label{thm:relative-finite-as-cover}
$S$ を Noetherian スキーム、$\mathcal V$ を階数四の局所自由層、
$\mathcal K\subset\mathcal V^{\otimes2}$ を階数七の局所直和因子とする。
二次環
\[
 \mathcal B=T_{\mathcal O_S}(\mathcal V)/(\mathcal K)
\]
について、各 $\mathcal B_n$ が有限局所自由で階数 $(n+1)^2$ を持ち、
全ての代数閉体からの射 $\bar s\to S$ に対して
$\mathcal B_{\bar s}$ が左右 noetherian な Koszul 環であると仮定する。
さらに $S$ の各アフィン開上で $\mathcal B$ は左右 noetherian とする。
この条件の任意の底変更への保存は仮定しない。

$G=\operatorname{Gr}_S(6,\mathcal K)$ とし、
普遍部分束を $\mathcal J$、商直線束を
$\mathcal L=\mathcal K/\mathcal J$ とする。
$G$ 上の二次環を
\[
 \mathcal A=T(\mathcal V)/\mathcal J,\qquad
 \mathcal D=\mathcal A^!=T(\mathcal V^*)/(\mathcal J^\perp)
\]
と書く。
次の条件を満たす局所閉部分スキーム $E_{\mathcal B}\subset G$ が存在し、
$S$ 上有限表示である。
\begin{enumerate}
\item $\mathcal A_3,\mathcal A_4$ が階数 $20,35$ の局所自由層である。
\item $\mathcal D_3,\mathcal D_4,\mathcal D_5$ がそれぞれ
階数 $4,1,0$ の局所自由層である。
次数零、一、二の階数は既に $1,4,6$ である。
\item 次数三の左右の積の像が、層として
\[
 \operatorname{im}(\mathcal L\otimes\mathcal V\to\mathcal A_3)
 =\operatorname{im}(\mathcal V\otimes\mathcal L\to\mathcal A_3)
\]
である。
\item $i=1,2$ に対する左右の積
\[
 \mathcal L\otimes\mathcal A_i\longrightarrow\mathcal A_{i+2},
 \qquad
 \mathcal A_i\otimes\mathcal L\longrightarrow\mathcal A_{i+2}
\]
が局所直和因子への単射である。
\item $i=1,2,3$ に対する積
$\mathcal D_i\otimes\mathcal D_{4-i}\to\mathcal D_4$
が完全な双線型組合せである。
\end{enumerate}

任意の Noetherian $S$-スキーム $T$ と
$\mathcal J_T\subset\mathcal K_T$ について、
$T\to G$ が $E_{\mathcal B}$ を経由することは、次と同値である。
\begin{enumerate}[label=(\alph*)]
\item $\mathcal A_T=T(\mathcal V_T)/(\mathcal J_T)$ の全次数が
有限局所自由であり、全ての幾何ファイバーが型 $(1,4,6,4,1)$ の
AS 正則環である。
\item $\mathcal L_T\subset(\mathcal A_T)_2$ は、局所生成元を選ぶと
左右正則正規元で生成される。
その正則性は任意の底変更後にも保たれ、
$\mathcal A_T/(\mathcal L_T)\simeq\mathcal B_T$ である。
\end{enumerate}
このとき $\operatorname{rank}(\mathcal A_T)_n=\binom{n+3}{3}$ であり、
左右の相対 Koszul 解消および双加群 Koszul 解消は長さ四で、
任意の底変更と両立する。
また、左右それぞれの相対 AS Gorenstein 条件は
\[
 \underline{\operatorname{Ext}}_{\mathcal A_T}^i
       (\mathcal O_T,\mathcal A_T)
 \simeq
 \begin{cases}
  (\mathcal D_T)_4\otimes\mathcal O_T(4),&i=4,\\
  0,&i\ne4
 \end{cases}
\]
となる。最高次数の直線束は局所的に自明化できる。
ここでの長さ四は相対解消の長さであり、
特異な基底上の環全体の大域次元を四と主張するものではない。

各アフィン開上で $\mathcal A_T$ が左右 noetherian であることは、
$\mathcal B_T$ が左右 noetherian であることと同値である。
従って $\mathcal B$ が全ての Noetherian 底変更の後にも左右 noetherian
であるという追加仮定の下では、$E_{\mathcal B}$ は
固定された $\mathcal B$ への全ての相対 noetherian AS 被覆を表現する。
この追加仮定を置かない場合は、$\mathcal B_T$ が左右 noetherian となる
テスト基底に限って同じ結論を用いる。
\end{theorem}
\begin{proof}
完全な証明は相対モジュライ論文 \cite[相対 AS 被覆の有限判定と fppf 像の反例]{QuadricModuli2026} に置く。
\end{proof}

\begin{remark}[相対被覆の存在として残る問題]
固定された平方根の相対軌道環 $\mathcal B$ に上の定理を適用できるとき、
被覆が fppf 局所に存在することは、有限表示の射
$E_{\mathcal B}\to S$ が fppf 局所に切断を持つことに等しい。
しかし、その射が滑らか・平坦・全射であるとは証明していない。
平方根そのものの相対存在と、得た被覆が相対 Serre 公式に適合することも
この定理の入力の外にある。
従ってこの有限表示への帰着だけを $\HHstack$ の代数性の証明としてはならない。
\end{remark}

\begin{example}[有限型射の fppf 像は代数的とは限らない]
\label{ex:fppf-image-nonalgebraic}
体 $k$ 上の有限型射
\[
 f:\mathbb A^2_{a,c}\longrightarrow\mathbb A^2_{a,b},
 \qquad (a,c)\longmapsto(a,ac)
\]
を考え、その fppf 層としての像を $I$ とする。
$I$ は $\mathbb A^2_{a,b}$ の部分層であるが、
代数空間では表現できない。
\end{example}
\begin{proof}
完全な証明は相対モジュライ論文 \cite[相対 AS 被覆の有限判定と fppf 像の反例]{QuadricModuli2026} に置く。
\end{proof}

この例は、幾何点の構成可能な像や scheme-theoretic image を求めることと、
非被約基底上の fppf 局所存在条件を表すこととの相違を示す。
\emph{この例から $\HHstack$ 自身の非代数性は結論しない。}
本研究では、平方根と被覆を忘れた部分関手の代数性を得るために、
さらに局所切断、相対認識、または同値関係の表示を証明する必要がある。
